\ProvidesFile{Global.tex}[Глобальные ресурсы]


\section{Глобальные ресурсы}
\label{section::Global}

Глобальные ресурсы в C++ является особым видом ресурсов.
Особенность их вытекает из того, что бывает несколько единиц трансляции, которые компилируются независимо.
И согласование глобальных ресурсов происходит на этапе линковки, когда все имена глобальные.
По этому нужно поговорить о том, как правильно работать с такими ресурсами, чтобы не огрести проблем.
Причина почему нам нужны глобальные ресурсы очень простая -- есть ситуации, когда без них просто нельзя.
Даже в языке есть свои глобальные ресурсы, например, потоки ввода вывода \verb"std::cout", \verb"std::cin", \verb"std::cerr".
Обратите внимание, что эти потоки гарантировано существуют всегда во время работы вашей программы и доступны из любой точки программы.
Часто на уровне операционной системы поддерживается какой-то один глобальный ресурс, например, единственный курсор мыши на всю ОС.
Вы можете иметь несколько контроллеров, но вы не можете для них иметь независимо несколько курсоров.
Самый популярный глобальный ресурс -- механизм логирования информации.

\subsection{Что хотим и в чем проблема}

Прежде всего перечислим, что может быть глобальным в C++:
\begin{enumerate}
\item Переменные

\item Функции

\item Классы

\item Шаблоны переменных

\item Шаблоны функций

\item Шаблоны классов
\end{enumerate}

\paragraph{Чего хотим}

Теперь рассмотрим конкретную задачу: мы хотим использовать Logger.
То есть мы хотим иметь возможность в любой точке программы сделать запись в некоторое заранее определенное хранилище информации для дальнейшего изучения.
В таком случае плохая идея передавать этот Logger в качестве аргумента в ваши функции, потому что вы замусорите интерфейс ваших функций и методов объектов нерелевантной информацией.
Например, если у вас есть функция для разложения целого числа на множители
\begin{cppcode}
std::vector<LongInt> factorize(const LonstInt&);
\end{cppcode}
то мы не хотели бы ее менять на
\begin{cppcode}
std::vector<LongInt> factorize(const LonstInt&, Logger&);
\end{cppcode}
Во-первых, это просто меняет api, что недопустимо, так как ломает весь код, который использует эту функцию.
Во-вторых, это ломает abi, теперь ваша функция имеет другое количество аргументов.
Но самое главное, в-третьих, Logger -- это логически лишний аргумент в методе \verb"factorize".
Ну потому что это не относится к той работе, которую делает функция.
Это какая-то работа по изучению поведения программы, отслеживанию проблем при ее исполнении.
У такого механизма скорее обслуживающий характер и должна быть возможность его быстро включить и выключить, не изменяя код нигде, кроме той точки, где выполняется логирование.

\paragraph{Наивное решение}

Наивное решение заключается в объявлении глобальной переменной для Logger-а и обращение к ней из разных единиц трансляции.
Давайте обсудим как это должно работать и почему этого механизма не достаточно.
Предположим у вас следующая структура файлов в проекте:%
{%
\ProvidesFile{GlobalCode.tex}[Код для раздела Глобальные ресурсы]


\newcommand{\codewidth}{5cm}

\newsavebox\globalh
\begin{lrbox}{\globalh}%
\begin{minipage}[\baselineskip]{\codewidth}
\texttt{Global.h}
\begin{cppcode}[numbers = none]
#pragma once

class Global {
public:
  void f() const;
private:
  ...
};
\end{cppcode}
\end{minipage}%
\end{lrbox}

\newsavebox\globalcpp
\begin{lrbox}{\globalcpp}%
\begin{minipage}[\baselineskip]{\codewidth}
\texttt{Global.cpp}
\begin{cppcode}[numbers = none]
#include "Global.h"
#include <iostream>

Global object;

void Global::f() const {
  std::cout << "f()\n";
}
\end{cppcode}
\end{minipage}%
\end{lrbox}

\newsavebox\maincpp
\begin{lrbox}{\maincpp}%
\begin{minipage}[\baselineskip]{\codewidth}
\texttt{main.cpp}
\begin{cppcode}[numbers = none]
#include "Global.h"

extern Global object;

int main() {
  object.f();
  return 0;
}
\end{cppcode}
\end{minipage}%
\end{lrbox}

\newsavebox\globalAh
\begin{lrbox}{\globalAh}%
\begin{minipage}[\baselineskip]{\codewidth}
\texttt{Global.h}
\begin{cppcode}[numbers = none]
#pragma once

class Global {
public:
  void f() const;
private
  ...
};

extern Global object;
\end{cppcode}
\end{minipage}%
\end{lrbox}

\newsavebox\mainAcpp
\begin{lrbox}{\mainAcpp}%
\begin{minipage}[\baselineskip]{\codewidth}
\texttt{main.cpp}
\begin{cppcode}[numbers = none]
#include "Global.h"



int main() {
  object.f();
  return 0;
}
\end{cppcode}
\end{minipage}%
\end{lrbox}

\newsavebox\globalBh
\begin{lrbox}{\globalBh}%
\begin{minipage}[\baselineskip]{\codewidth}
\texttt{A.h}
\begin{cppcode}[numbers = none]
#pragma once

class A {
public:
  A();
private
  ...
};

extern A a;
\end{cppcode}
\end{minipage}%
\end{lrbox}

\newsavebox\globalBcpp
\begin{lrbox}{\globalBcpp}%
\begin{minipage}[\baselineskip]{\codewidth}
\texttt{A.cpp}
\begin{cppcode}[numbers = none]
#include "A.h"
#include "Global.h"

A a;

A::A() {
  object.f();
}
\end{cppcode}
\end{minipage}%
\end{lrbox}

\newsavebox\mainBcpp
\begin{lrbox}{\mainBcpp}%
\begin{minipage}[\baselineskip]{\codewidth}
\texttt{main.cpp}
\begin{cppcode}[numbers = none]
#include "Global.h"
#include "A.h"


int main() {
  object.f();
  return 0;
}
\end{cppcode}
\end{minipage}%
\end{lrbox}


\[
\xymatrix{
  *+[F:<3pt>]{\usebox{\globalh}}&{}\\
  *+[F:<3pt>]{\usebox{\globalcpp}}\ar[u]&*+[F:<3pt>]{\usebox{\maincpp}}\ar[ul]\\
}
\]
Давайте объясним что тут происходит.
В начале разберемся с тем, что происходит при компиляции проекта:
\begin{enumerate}
\item В проекте две единицы трансляции:
\begin{enumerate}
\item \verb"Global.cpp"

\item \verb"main.cpp"
\end{enumerate}

\item При компиляции \verb"Global.cpp" компилятор видит код:
\begin{cppcode}
// Global TU

class Global {
public:
  void f() const;
private:
  ...
};

Global object;

void Global::f() const {
  std::cout << "f()\n";
}
\end{cppcode}

В данной единице трансляции он видит два объекта:
\begin{enumerate}
\item Глобальная переменная \verb"object" типа \verb"Global".

\item Метод \verb"Global::f".
\end{enumerate}
Значит компилятор создает бинарный файл, в который складывает глобальную переменную \verb"object" и код для функции \verb"Global::f".
И получается логически следующее:
\begin{cppcode}
// Global.bin

Global object

void Global::f:
  std::cout << "f()\n";
\end{cppcode}
Обратите внимание, что определение класса нужно лишь для того, чтобы суметь скомпилировать код.
Оно необходимо, чтобы понять сколько памяти выделять под \verb"object" при его инициализации в бинарном файле.%
\footnote{Тут на самом деле еще неявно генерируются дефолтный конструктор и деструктор для \texttt{object}.}

\item При компиляции \verb"main.cpp" компилятор видит код:
\begin{cppcode}
// Global main

class Global {
public:
  void f() const;
private:
  ...
};

extern Global object;

int main() {
  object.f();
  return 0;
}
\end{cppcode}

\item При компиляции этой единицы трансляции, компилятор видит три объекта:
\begin{enumerate}
\item Глобальная переменная \verb"object", помеченная специальным словом \verb"extern"

\item Функция \verb"Global::f", у которой отсутствует определение.

\item Функция \verb"main" -- точка входа в программу.
\end{enumerate}
При компиляции это единицы трансляции происходит следующее:
\begin{enumerate}
\item В бинарный файл складывается точка входа в программу -- код функции \verb"main".

\item Переменная помеченная словом \verb"extern" не создается в бинарном файле.
Это объявление нужно компилятору в строке~13, чтобы понять, что \verb"object.f()" -- это вызов метода \verb"f" у объекта \verb"object" класса \verb"Global".
Но мы говорим компилятору, что настоящая переменная должна быть заведена в памяти не в текущем бинарнике, а в каком-то другом и пусть линковщик ее сам найдет.
А если не найдет, то у вас будут проблемы.

\item После чего компилятор видит строчку~14, где находится вызов метода класса.
Это мы понимаем из объявления класса \verb"Global", который был вставлен директивой \verb"include".
И так как для нее нет определения, то мы просто ставим точку вызова этой функции в бинарник.
\end{enumerate}
В результате получится следующий бинарный файл:
\begin{cppcode}
// main.bin

main:
  call object.f()
\end{cppcode}

\item Теперь в работу включается линковщик, который сливает вместе бинарные файлы так:
\begin{cppcode}
// main.exe

Global object // create in static memory

main:
  call object.f()
  
Global::f:
  std::cout << "f()\n";
\end{cppcode}
\end{enumerate}
Таким образом про \verb"extern" можно думать, как про ссылку на глобальный объект, которая как бы делается на уровне линковщика.

\paragraph{Проблемы}

Основная идея использования \verb"h" файлов в модели компиляции C++ в том, чтобы зафиксировать в них имена функций, классов и их методов, чтобы при изменении названия метода, компилятор мог нам подсказать в чем проблема и найти все точки в коде, где нужно внести изменения в код.
В версии выше глобальная переменная, имя которой должно быть точкой доступа к ресурсу, объявляется и определяется внутри \verb"cpp" файла, что скрывает ее изменение ее имени от пользователя.
Однако, эту проблему можно легко исправить так:
\[
\xymatrix{
  *+[F:<3pt>]{\usebox{\globalAh}}&{}\\
  *+[F:<3pt>]{\usebox{\globalcpp}}\ar[u]&*+[F:<3pt>]{\usebox{\mainAcpp}}\ar[ul]\\
}
\]
В таком виде это код уже более или менее удобен для использования, но остается важная техническая проблема, которую я хочу обсудить ниже.

\paragraph{Порядок инициализации статических переменных}

Предположим, что у нас теперь есть два глобальных объекта:
\[
\xymatrix{
  *+[F:<3pt>]{\usebox{\globalAh}}&{}&*+[F:<3pt>]{\usebox{\globalBh}}\\
  *+[F:<3pt>]{\usebox{\globalcpp}}\ar[u]&*+[F:<3pt>]{\usebox{\mainAcpp}}\ar[ul]\ar[ur]&*+[F:<3pt>]{\usebox{\globalBcpp}}\ar[u]\ar@/_20pt/[ull]\\
}
\]
Конструкторы глобальных объектов отрабатывают до старта функции \verb"main".
Внутри одной единицы трансляции глобальные объекты инициализируются в порядке их определения.
Однако, между разными единицами трансляции порядок инициализации никоим образом не специфицирован.
А это значит, что \verb"a" может быть сконструирован раньше, чем \verb"object".
Но тогда в конструкторе \verb"A::A()" будет обращение к еще не инициализированному объекту в памяти \verb"object".
Вот эта трудность является основной проблемой в работе с глобальными объектами -- нет гарантии, что глобальный объект создан к моменту его использования.
Если вы используете глобальные объекты после старта функции \verb"main" и до выхода из нее, то такой проблемы нет.
Однако, если же вы используете глобальные объекты в деструкторах глобальных объектов, то проблема может всплыть.
}

\subsection{Singleton}
\label{section::Singleton}

\paragraph{Scott Mayers Singleton}

Давайте начнем с классического решения -- Scott Mayers Singleton.
Основная проблема для нас -- глобальные переменные не обязаны быть инициализированы к моменту обращения к ним.
Эта проблема решается специальным видом глобальных переменных -- статические переменные внутри функции.
Давайте посмотрим на следующий кусок кода:
\begin{cppcode}
class A { ... };

A& instance_A() {
  static A a;
  return a;
}
\end{cppcode}
В этом случае функция \verb"access" возвращает ссылку на статическую переменную \verb"a".
В этом случае \verb"a" является глобальной переменной, которая доступна только внутри функции \verb"access".
Но у этой переменной особое поведение при инициализации.
По стандарту эта переменная инициализируется при первом обращении к функции \verb"access".
Более того, гарантируется, что инициализация такой статической переменной является потокобезопасной.
Что делает это решение максимально простым и надежным в использовании.

Использование такого глобального ресурса выглядит так:
\begin{cppcode}
int main() {
  instance_A().f();
  return 0;
}
\end{cppcode}

\paragraph{Шаблонное решение}

Теперь давайте завернем Scott Mayers Singleton в шаблонное решение.
Можно написать следующее решение:
\begin{cppcode}
template<class T, class Tag>
class Global {
public:
  using type = T;
  using tag = Tag;

  static T& instance() {
    static T data;
    return data;
  }
};
\end{cppcode}
В этом случае можно использовать это решение так:
\begin{cppcode}
class A {
public:
  void f() const;
  int g(int) const;
};

using GAmain = Global<A, struct Main>;
using GAopt = Global<A, struct Optional>;

int main() {
  GAmain::instance().f();
  int x = GAmain::instance().g(1);
  
  GAopt::instance().f();
  return 0;
}
\end{cppcode}
В этом случае у шаблона два параметра:
\begin{enumerate}
\item Тип \verb"T" глобального объекта, к которому мы получаем доступ.

\item Вспомогательный тип \verb"Tag".
Этот тип нужен, чтобы суметь создать разные глобальные объекты одного типа.
Действительно, если бы у нас был бы только один параметр, то мы могли бы написать лишь
\begin{cppcode}
using GlobalA = Global<A>;
\end{cppcode}
В этом случае для класса типа \verb"A" мы больше не сможем создать других глобальных объектов этого же типа.
Тег по сути позволяет создавать разные объекты одного типа.
А заодно имя тега подсказывает назначение глобального объекта.
\end{enumerate}

\paragraph{Изменение синтаксиса}

Обратите внимание, как нелепо выглядит обращение к методу глобального объекта
\begin{cppcode}
GlobalA::access().f();
\end{cppcode}
Можно добиться более приятного и удобного синтаксиса с помощью прокси объектов аксессоров.
Сделать это можно так:
\begin{cppcode}
template<class Global>
class Accessor {
public:
  Global::type* operator->() {
    return &Global::instance();
  }
  const Global::type* operator->() const {
    return &Global::instance();
  }
};
\end{cppcode}
То есть объект аксессор получает доступ к глобальному ресурсу через оператор стрелку.
Обратите внимание, что \verb"const" проносится сквозь стрелку.
Это связано с тем, что константный аксессор должен давать только доступ на чтение, а не константный -- любой.
Тогда использование будет выглядеть так:
\begin{cppcode}
using GlobalA = Global<A, struct Main>;
using AccessA = Accessor<GlobalA>;

int main() {
  AccessA a;
  a->f();
  int x = a->g(1);
  return 0;
}
\end{cppcode}
Так код становится более читаемым и менее замусоренным.

\subsection{Singleton с параметрами}
\label{section::SingletonParams}

Следующий большой вопрос: а что делать, если нам не подходит дефолтный конструктор для глобального объекта?
Как тогда организовать код, чтобы можно было создать этот объект с нужными параметрами?
Тут есть следующая проблема: только первое обращение к глобальному ресурсу должно его инициализировать, а все остальные обращения должны его просто использовать.
Это автоматически означает, что мы должны знать, какое обращение является первым и организовать все остальные обращения после инициализации.
В частности это означает, что такие глобальные ресурсы нужно инициализировать руками до первого использования.
Для простейшего решения нам нужна глобальная переменная, к которой есть доступ из двух вспомогательных функций: одна инициализирует эту переменную, а другая предоставляет доступ.
По умолчанию в языке нет механизма, чтобы шарить между функциями одну вспомогательную глобальную переменную.
Потому нужны другие механизмы для ограничения несанкционированного доступа к ресурсу.
Кроме того, нам теперь нужно предоставлять доступ к неициализированному ресурсу.

\paragraph{Концепция решения}

Ресурс можно хранить на куче или в статической памяти.
\begin{center}
\begin{tabular}{cc}
{
\begin{minipage}[\baselineskip]{8cm}
В статической памяти.
\begin{cppcode}[numbers=none]
class A { ... };

std::optional<A>& data() {
  static std::optional<A> a;
  return a;
}
\end{cppcode}
\end{minipage}
}&{
\begin{minipage}[\baselineskip]{8cm}
На куче.
\begin{cppcode}[numbers=none]
class A { ... };

std::unique_ptr<A>& data() {
  static std::unique_ptr<A> a;
  return a;
}
\end{cppcode}
\end{minipage}
}\\
\end{tabular}
\end{center}
Теперь можно создать методы инициализации и доступа к данным.
\begin{center}
\begin{tabular}{cc}
{
\begin{minipage}[\baselineskip]{8cm}
В статической памяти.
\begin{cppcode}[numbers=none]
template<class... Args>
A& initialize(Args&&... args) {
  assert(!data());
  data() = std::make_optional<A>(
             std::forward<Args>(args)...);
  return *data();
}

A& instance() {
  assert(data());
  return *data();
}
\end{cppcode}
\end{minipage}
}&{
\begin{minipage}[\baselineskip]{8cm}
На куче.
\begin{cppcode}[numbers=none]
template<class... Args>
A& initialize(Args&&... args) {
  assert(!data());
  data() = std::make_unique<A>(
             std::forward<Args>(args)...);
  return *data();
}

A& instance() {
  assert(data());
  return *data();
}
\end{cppcode}
\end{minipage}
}\\
\end{tabular}
\end{center}
Тогда можно использовать данный глобальный ресурс следующим образом:
\begin{cppcode}
int main() {
  //instance().f(); // assert(data()) failed! the object is not initialized yet

  initialize(1, "123");
  //initialize(1, "abc"); // assert(!data()) failed! reinitialization of the object
  
  instance().f();
  instance().f();
  return 0;
}
\end{cppcode}

\paragraph{Шаблонное решение}

Теперь завернем это в шаблонное решение.
Я напишу на примере \verb"optional", для указателя делается аналогично.
\begin{cppcode}
template<class T, class Tag>
class Global {
public:
  using type = T;
  using tag = Tag;

  template<class... Args>
  static T& initialize(Args&&... args) {
    assert(!data());
    data() = std::make_optional<T>(std::forward<Args>(args)...);
    return *data();
  }
  static T& instance() {
    assert(data());
    return *data();
  }
private:
  static std::optional<T>& data() {
    static std::optional<T> object;
    return object;
  }
};
\end{cppcode}
Так же можно сделать класс аксессор для удобства:
\begin{cppcode}
template<class G>
class Accessor {
public:
  using T = typename G::type;

  const T* operator->() const {
    return &G::instance();
  }
  T* operator->() {
    return &G::instance();
  }
};
\end{cppcode}
И теперь этим можно пользоваться так:
\begin{cppcode}
class A { ... };

using GlobalA = Global<A, struct Main>;
using AccessA = Accessor<GlobalA>;

int main() {
  GlobalA::initialize(1, "123");
  GlobalA::instance().f();
  
  AccessA a;
  a->f();
  return 0;  
}
\end{cppcode}
При наличии класса аксессора можно закрыть доступ по \verb"instance", чтобы избежать уродливого синтаксиса так:
\begin{center}
\begin{tabular}{cc}
{
\begin{minipage}[\baselineskip]{8cm}
\begin{cppcode}[numbers=none]
template<class T, class Tag>
class Global {
public:
  using type = T;
  using tag = Tag;

  template<class... Args>
  static T& initialize(Args&&... args) {
    assert(!data());
    data() = std::make_optional<T>(
               std::forward<Args>(args)...);
    return *data();
  }
protected:
  static T& instance() {
    assert(data());
    return *data();
  }
private:
  static std::optional<T>& data() {
    static std::optional<T> object;
    return object;
  }
};
\end{cppcode}
\end{minipage}
}&{
\begin{minipage}[\baselineskip]{8cm}
\begin{cppcode}[numbers=none]
template<class G>
class Accessor : G {
public:
  using T = typename G::type;

  const T* operator->() const {
    return &G::instance();
  }
  T* operator->() {
    return &G::instance();
  }
};

class A { ... };

using GlobalA = Global<A, struct Main>;
using AccessA = Accessor<GlobalA>;

int main() {
  GlobalA::initialize(1, "123");
  AccessA a;
  a->f();
  return 0;  
}
\end{cppcode}
\end{minipage}
}\\
\end{tabular}
\end{center}
В этом случае \verb"Accessor" является производным классом, что позволяет ему иметь доступ к защищенным методам.
При этом это никак не влияет на содержимое объекта класса \verb"Accessor", потому что базовый класс с глобальным ресурсом не содержит никаких переменных и методов класса (только статические = глобальные методы).
Так же обратите внимание на проброску \verb"const" сквозь оператор стрелку.

\subsection{Классическое решение}

Если вы посмотрите в различную литературу по паттернам проектирования, то увидите, что авторы поголовно предлагают следующую имплементацию Singleton:
\begin{cppcode}
class A {
public:
  static A& instance() {
    static A a;
    return a;
  }
  
  void f() const {
    std::cout << "A::f()\n";
  }
  
  int g(int x) const {
    return 2 * x;
  }
private:
  A() = default;
  A(const A&) = default;
  A(A&&) noexcept = default;
  A& operator=(const A&) = default;
  A& operator=(A&&) noexcept = default;
  ~A() = default;
};
\end{cppcode}
И тогда можно пользоваться классом следующим образом:
\begin{cppcode}
int main() {
  A::instance().f();
  std::cout << "value = " << A::instance().g(1) << '\n';
  return 0;
}
\end{cppcode}
У такого подхода есть несколько недостатков:
\begin{enumerate}
\item Данный класс выполняет две разные задачи, не связанные друг с другом:
\begin{enumerate}
\item Обеспечить функциональность класса.
Это имплементации методов \verb"A::f" и \verb"A::g".

\item Обеспечение глобального доступа к единственному объекту класса \verb"A".
Это имплементация метода \verb"instance".
\end{enumerate}

\item Нельзя переиспользовать код для работы с единственным глобальным объектом.

\item При таком подходе на уровне языка запрещено несколько объектов класса \verb"A".
Это не всегда бывает нужно.
\end{enumerate}
Есть еще одна психолоически неприятная вещь, которая возникает при таком подходе к Singleton-ам, -- такой код скрывает зависимости между разными Singleton-ами.
Ниже будет написано, что в точности я под этим подразумеваю, и в чем может быть проблема.

\subsection{Зависимые Singleton-ы}

Предположим, что у нас есть два глобальных ресурса \verb"A" и \verb"B" и ресурс \verb"B" требует для корректной работы ресурс \verb"A".
Существует два принципиальных способа как реализовать такую зависимость.
Первый подход -- напрямую обращаться к глобальному ресурсу.
Я специально сделаю тут классическую имплементацию.
\begin{center}
\begin{tabular}{cc}
{
\begin{minipage}[\baselineskip]{8cm}
\begin{cppcode}[numbers=none]
class A {
public:
  static A& instance();
  
  void f() const {
    std::cout << "A::f()\n";
  }
  
private:
  A() = default;
  A(const A&) = default;
  A(A&&) noexcept = default;
  A& operator=(const A&) = default;
  A& operator=(A&&) noexcept = default;
  ~A() = default;  
};
\end{cppcode}
\end{minipage}
}&{
\begin{minipage}[\baselineskip]{8cm}
\begin{cppcode}[numbers=none]
class B {
public:
  static B& instance();
  
  void f() const {
    A::instance().f();
    std::cout << "B::f()\n";
  }
private:
  B() = default;
  B(const B&) = default;
  B(B&&) noexcept = default;
  B& operator=(const B&) = default;
  B& operator=(B&&) noexcept = default;
  ~B() = default;  
};
\end{cppcode}
\end{minipage}
}\\
\end{tabular}
\end{center}
Обратите внимание, что так как на уровне языка у нас класс \verb"A" имеет только один объект возвращаемый \verb"A::instance", то нам ничего не мешает логически без зазрения совести напрямую в методах \verb"B" обращаться к методам этой переменной.
В этом случае код не выглядит подозрительным и даже нормально работает.
Теперь давайте сделаем то же самое с помощью шаблонного Singleton описанного в разделе~\ref{section::Singleton}.
\begin{center}
\begin{tabular}{cc}
{
\begin{minipage}[\baselineskip]{8cm}
\begin{cppcode}[numbers=none]
class A {
public:

  void f() const {
    std::cout << "A::f()\n";
  }
};


using GlobalA = Global<A, struct Main>;
using AccessA = Accessor<GlobalA>;
\end{cppcode}
\end{minipage}
}&{
\begin{minipage}[\baselineskip]{8cm}
\begin{cppcode}[numbers=none]
class B {
public:
  void f() const {
    AssessA a;  // this is strange
    a->f();
    std::cout << "B::f()\n";
  }
};

using GlobalB = Global<B, struct Main>;
using AccessB = Accessor<GlobalB>;
\end{cppcode}
\end{minipage}
}\\
\end{tabular}
\end{center}
Обратите внимание, что теперь, когда синтаксически класс \verb"A" не обязан иметь один единственный глобальный объект, то обращение к методу\verb"A::f" внутри класса \verb"B" выглядит странным.
Теперь классу \verb"B" не достаточно знать только класс \verb"A", а надо еще знать какой именно глобальный ресурс этого типа можно использовать.
Более того, возникает вопрос: а что если я создам переменную типа \verb"B" локально (или создам несколько локальных переменных), то это нормально что они все ссылаются на один и тот же фиксированный глобальный объект?
При использовании в виде глобальных ресурсов действительно разницы нет, но второй подход к написанию кода помогает увидеть, что у нас есть неявная зависимость, которую лучше сделать явной.
И видим мы это потому, что наш класс теперь используется не только в виде глобального ресурса.

В этом случае можно явно сохранить в \verb"B" указатель на требуемый ресурс.
\begin{cppcode}
class B {
public:
  B(A& a) : a_(&a) {}
  void f() const {
    a_->f();
    std::cout << "B::f()\n";
  }
private:
  A* a_;
};
\end{cppcode}
Теперь зависимость от класса \verb"A" явно выражена в аргументах конструктора.
Однако, мы не можем теперь использовать \verb"B" как глобальный ресурс с дефолтным конструктором, так как у нас больше его нет.
В этом случае можно добавить дефолтный конструктор, в котором инициализация идет с помощью глобального ресурса \verb"A".
\begin{cppcode}
class B {
public:
  B() = default;
  B(A& a) : a_(&a) {}
  void f() const {
    a_->f();
    std::cout << "B::f()\n";
  }
private:
  A* a_ = &A::instance(); // need to implement this method in Global
};
\end{cppcode}
По хорошему порядок создания и разрушения глобальных объектов должен гарантировать, что объект \verb"A::instance" доступный по указателю \verb"a_" будет жив, пока жив объект типа \verb"B".


Другой подход -- сделать инициализацию всех объектов явной и передавать все аргументы явно.
Тут я воспользуюсь шаблоном \verb"Global" из раздела~\ref{section::SingletonParams}.
\begin{center}
\begin{tabular}{cc}
{
\begin{minipage}[\baselineskip]{8cm}
\begin{cppcode}[numbers=none]
class A {
public:

  void f() const {
    std::cout << "A::f()\n";
  }
};




using GlobalA = Global<A, struct Main>;
using AccessA = Accessor<GlobalA>;
\end{cppcode}
\end{minipage}
}&{
\begin{minipage}[\baselineskip]{8cm}
\begin{cppcode}[numbers=none]
class B {
public:
  B(A& a) : a_(&a) {}
  void f() const {
    a_->f();
    std::cout << "B::f()\n";
  }
private:
  A* a_;
};

using GlobalB = Global<B, struct Main>;
using AccessB = Accessor<GlobalB>;
\end{cppcode}
\end{minipage}
}\\
\end{tabular}
\end{center}
Тогда это будет использоваться так
\begin{cppcode}
int main() {
  GlobalA::initialize();
  GlobalB::initialize(*AccessA());
  
  AccessA a;
  a->f();
  AccessB b;
  b->f();
  return 0;
}
\end{cppcode}
Здесь я подразумеваю, что у аксессора кроме оператора стрелки есть еще оператор \verb"operator*", который возвращает ссылку на глобальный ресурс.
Таким образом мы получим глобальные ресурсы иницализированные в правильном порядке.
Отсутствие явной зависимости между классами \verb"A" и \verb"B" в аргументах конструктора для \verb"B" влечет к проблемам при инициализации.
Действительно, если написать:
\begin{center}
\begin{tabular}{cc}
{
\begin{minipage}[\baselineskip]{8cm}
\begin{cppcode}[numbers=none]
class A {
public:

  void f() const {
    std::cout << "A::f()\n";
  }
};


using GlobalA = Global<A, struct Main>;
using AccessA = Accessor<GlobalA>;
\end{cppcode}
\end{minipage}
}&{
\begin{minipage}[\baselineskip]{8cm}
\begin{cppcode}[numbers=none]
class B {
public:
  void f() const {
    AccessA a;
    a->f();
    std::cout << "B::f()\n";
  }
};

using GlobalB = Global<B, struct Main>;
using AccessB = Accessor<GlobalB>;
\end{cppcode}
\end{minipage}
}\\
\end{tabular}
\end{center}
То при использовании возможны следующие две ситуации:
\begin{center}
\begin{tabular}{cc}
{
\begin{minipage}[\baselineskip]{8cm}
\begin{cppcode}[numbers=none]
int main() {
  GlobalA::initialize();
  GlobalB::initialize();
  
  AccessA a;
  a->f();
  AccessB b;
  b->f();
  return 0;
}
\end{cppcode}
\end{minipage}
}&{
\begin{minipage}[\baselineskip]{8cm}
\begin{cppcode}[numbers=none]
int main() {
  GlobalB::initialize();  // UB
  GlobalA::initialize();
  
  AccessA a;
  a->f();
  AccessB b;
  b->f();
  return 0;
}
\end{cppcode}
\end{minipage}
}\\
\end{tabular}
\end{center}
В этом случае ничто не мешает нам вызвать инициализацию в неверном порядке.
А это означает, что при обращении в конструкторе \verb"B" к объекту \verb"A", мы получим обращение к еще не созданному объекту.
А это undefined behaviour.

\subsection{Преинициализация ресурсов}

Существует еще одна ситуация, когда нам может понадобиться вызвать инициализацию глобальных ресурсов руками.
Обратите внимание, что для Scott Mayers Singleton первое обращение к методу \verb"instance" дорогое, потому что требует инициализации ресурсов.
Если для вас это критично и вы хотите заранее инициализировать все глобальные ресурсы, то можно заранее вызвать метод обращения к ресурсу. 
Однако, если вы используете шаблон из раздела~\ref{section::Singleton}, то на стороне пользователя код будет выглядет странно и не понятно:
\begin{cppcode}
class A { ... };

using GlobalA = Global<A, struct Main>;
using AccessA = Accessor<GlobalA>;

int main() {
  GlobalA::instance(); // initialization
  
  GlobalA::instance().f();
  AccessA a;
  a->f();
  return 0;
}
\end{cppcode}
В этом случае лучше перейти на решение из раздела~\ref{section::SingletonParams} даже при использовании дефолтных конструкторов.
\begin{cppcode}
class A { ... };

using GlobalA = Global<A, struct Main>;
using AccessA = Accessor<GlobalA>;

int main() {
  GlobalA::initialize();
  
  GlobalA::instance().f();
  AccessA a;
  a->f();
  return 0;
}
\end{cppcode}
В этом случае код говорит в точности то, что он делает.

\subsection{Рефакторинг Singleton-ов}
\label{section::SingletonRefactoring}

Если вы столкнулись с кодом, в котором используются Singleton-ы и вы хотите от них избавиться, не ломая при этом уже написанный код, то есть одна разумная тактика, которую я хочу продемонстрировать.

\paragraph{Функции}

Предположим, что у вас есть функция, которая внутри использует глобальный ресурс.
Пусть для определенности будет что-то в таком духе:
\begin{cppcode}
int function(int x, double y, const B& z) {
  // some code
  GlobalA::instance().f();
  //some code
  return /*something*/;
}
\end{cppcode}
Теперь мы хотим избавиться от этой зависимости от глобального объекта.
Его для этого его можно передать явно в качестве аргумента:
\begin{cppcode}
int function(int x, double y, const B& z, A& obj) {
  // some code
  obj.f();
  //some code
  return /*something*/;
}
\end{cppcode}
Однако. таким образом мы сломали api, а значит и код, который использует нашу функцию.
\begin{cppcode}
int main() {
  B b;
  int z = func(1, 2.0, b); //  compilation error
  return 0;
}
\end{cppcode}
Теперь такой вызов перестанет компилироваться.
Исправить ситуацию можно так:
\begin{cppcode}
int function(int x, double y, const B& z, A& obj =  GlobalA::instance()) {
  // some code
  obj.f();
  //some code
  return /*something*/;
}
\end{cppcode}
Так мы решили проблему и код
\begin{cppcode}
int main() {
  B b;
  int z = func(1, 2.0, b); // compiles
  return 0;
}
\end{cppcode}
Однако, мы таким образом сломали abi.
Действительно до этого была функция $3$ аргументов, а теперь функция $4$ аргументов.
А это другая работа со стеком вызовов.
Если это принципиально, то можно воспользоваться перегрузкой функций:
\begin{cppcode}
int function(int x, double y, const B& z, A& obj) {
  // some code
  obj.f();
  //some code
  return /*something*/;
}

int function(int x, double y, const B& z) {
  return function(x, y, z, GlobalA::instance());
}
\end{cppcode}
Теперь старая функция определена в терминах новой функции.
И старый код будет работать
\begin{cppcode}
int main() {
  B b;
  int z = func(1, 2.0, b); // compiles
  return 0;
}
\end{cppcode}
И abi при вызове функции сохранится.
Изменится имплементация, но это ожидаемое поведение.
Подобная техника позволяет оставить там, где не надо ничего менять, код, который работает с глобальными ресурсами.
А там, где нужна работа с чем-то локальным, использовать перегруженную версию с явным указанием объекта.

\paragraph{Методы класса и объекты класса}

Теперь предположим, что у вас есть класс, который внутри себя в своих методах использует глобальный ресурс.
Например что-то такое:
\begin{cppcode}
class A {
public:
  void function() const {
    // some code
    GlobalB::instance().f();
    // some code
  }
private:
  S data_;
};
\end{cppcode}
В этом случае есть несколько вариантов:
\begin{enumerate}
\item Передавать явно в аргументах методов.

\item Сделать глобальным ресурсом внутри класса.

\item Сделать локальную ссылку внутри каждого объекта.

\item Получать ссылку на нужный объект через Observer порт.
\end{enumerate}
В первом случае это выглядит так:
\begin{cppcode}
class A {
public:
  void function(GlobalB& obj) const {
    // some code
    obj.f();
    // some code
  }
  void function() const {
    function(GlobalB::instance());
  }
private:
  S data_;
};
\end{cppcode}
Второй подход подразумевает, что мы для класса \verb"A" создадим глобальный ресурс спрятанный внутри самого класса так:
\begin{cppcode}
class A {
  using MyGlobalB = Global<B, Main>;
public:
  void function() const {
    // some code
    MyGlobalB::instance().f(); // A::MyGlobalB is used here
    // some code
  }
private:
  S data_;
};
\end{cppcode}
Обратите внимание, что полное имя тега будет \verb"A::Main", а потому у нас точно не будет случайного пересечения с тегом любого другого глобального ресурса.

В третьем подходе, мы сохраняем версию глобального объекта внутри при конструировании:
\begin{cppcode}
class A {
public:
  A() = default;
  A(B& obj) : obj_(obj) {}
  void function() const {
    // some code
    obj().f();
    // some code
  }

private:
  B&  obj() const {
    return obj_.get();
  }
  
  S data_;
  std::reference_wrapper<B> obj_ = GlobalB::instance();
};
\end{cppcode}
Тут хранение происходит в \verb"reference_wrapper" для обеспечение value semantics класса.
Можно было бы использовать указатель, но тогда за инвариант, что указатель не нулевой должен был бы отвечать класс.
Вместо этого за этот инвариант отвечает само поле класса.
Важно, чтобы ссылка на ресурс была жива все время жизни объекта класса \verb"A".
Ссылка на глобальный ресурс в этом случае работает безопасно.
Так же чтобы избежать длинных вызовов вида
\begin{cppcode}[numbers = none]
obj_.get().f();
\end{cppcode}
вводится вспомогательная функция для обращения к ссылке хранящейся внутри \verb"reference_wrapper".
Такой подход меняет abi, потому что меняется layout для объекта класса, как минимум его размер увеличивается на один указатель.
Зато теперь объект класса служит как контекст для передачи ресурса своим полям.
Например можно протолкнуть зависимость так:
\begin{cppcode}
class A {
public:
  A() = default;
  A(B& obj) : obj_(obj) {}
  ...
private:
  B&  obj() const {
    return obj_.get();
  }

  std::reference_wrapper<B> obj_ = GlobalB::instance();
  S data_ = obj();
};
\end{cppcode}
Обратите внимание, что для такой махинации нам надо будет переставить поля, чтобы инициализация прошла в правильном порядке.

И теперь последний четвертый подход.
В этом случае у нас будут два класса, которые выступают в роли модулей приложения: \verb"A" и \verb"C" (который и будет оповещать нужной версией объекта класса \verb"B").
Класс \verb"A" модифицируется так:
\begin{cppcode}
class A {
public:
  void function() const {
    // some code
    if (!port_.has_data()) // Need to support this case probably with different behavior
      return;
    port_->f();
    // some code
  }
  Observer<B>* port() {
    return &port_;
  }
private:
  Observer<B> port_;
  S data_;
};
\end{cppcode}
Класс \verb"C" будет выглядеть так:
\begin{cppcode}
class C {
public:
  void subscribe(Observer<B>* port) {
    port_.subscribe(port);
  }
  void set(const B& b) {
    b_ = b;
    port_.notify();
  }
private:
  B b_;
  Observable<B> port_ = /*binding to b_*/;
};
\end{cppcode}
И теперь для использования надо лишь соединить объекты:
\begin{cppcode}
A a;
C c;
c.subscribe(a.port());
c.set(B{});
c.set(B{});
\end{cppcode}
Таким образом удается избавиться от неявной зависимости.
Единственный минус такого подхода в том, что приходится во всех методах предусматривать случай отсутствия у \verb"A" объекта типа \verb"B", потому что порт мог быть не подсоединен.

\subsection{Logging}
\label{section::Logging}

Предположим, что мы хотим использовать удобный механизм логирования информации.
Для этого нам нужен объект сериализатор%
\footnote{Про сериализатор можно чуть-чуть почитать тут~\ref{section::Serialization}.}
и шаблон для получения доступа к глобальному ресурсу.
Выглядеть это может так
\begin{cppcode}
class Logger {
public:
   template<class T>
   Logger& operator<<(const T&) {
     /*Implemetation for each class T*/
   }
   ...
private:
  ...
};

inline Logger& main_logger() {
  using MainLogger = Global<Logger, struct Main>;
  return MainLogger::instance();
}
\end{cppcode}
Теперь использование будет выглядеть как-то так
\begin{cppcode}
void function(int x, const A& y) {
  // some code
  main_logger() << "x = " << x << ", y = " << y << '\n';
  // some code
}
\end{cppcode}

\subsection{Подсчет количества операций}
\label{section::OpCounting}

Предположим, вы хотите понять на сколько аккуратно вы пользуетесь value semantics.
Например, вы не делаете лишних копий там, где этого не нужно.
Для этого можно организовать подсчет количества копирований объектов при исполнении.
Общая идея решения такая: создадим класс, который умеет считать количество вызовов своих конструкторов, деструкторов и перемещающих операторов, а потом унаследуемся от этого класса.
Подсчет можно делать глобально на протяжении жизни всей программы, а можно локально только в рамках какого-то scope.
Сразу оговорюсь, что решение ниже не поддерживает многопоточность.
Многопоточное решение -- это совсем другой разговор, который я тут вести не хочу.
Так же отмечу, что хоть мы и считаем только конструкторы, присваивания и деструкторы, но вообще говоря подобным методом можно считать количество вызовов любого метода или операций.
Например так можно автоматизировать подсчет количества проделанных операций над объектами.
Это позволяет при запуске алгоритма проверить реальное количество проделанных операций с более или менее локальным вмешательством в код.

\paragraph{Использование}

Прежде чем работать над имплементацией, давайте обсудим, как мы видим использование подобного функционала.
Вот пусть у нас есть класс \verb"A" и мы хотим посчитать во время работы программы сколько раз срабатывают операции копирования в программе.
В этом случае добавим публичное наследование от служебного класса:
\begin{cppcode}
class A : public CTorCheck<A> {
public:
  A() = default;
  A(int);
  ...
private:
  ...
};
\end{cppcode}
Все больше ничего  делать не надо.
Теперь при исполнении следующего кода:
\begin{cppcode}
int main() {
  algorithm(); // some algorithm using objects of class A
  return 0;
}
\end{cppcode}
после выхода из функции \verb"main", мы получим информацию о количестве выполненных операций:
\begin{enumerate}
\item Количество вызванных конструкторов.

\item Количество вызванных копирующих конструкторов.

\item Количество вызванных перемещающих конструкторов.

\item Количество вызванных копирующих присваиваний.

\item Количество вызванных перемещающих присваиваний.

\item Количество вызванных деструкторов.
\end{enumerate}

\paragraph{Имплементация считающего объекта}

Давайте в начале создадим класс \verb"Counter", который будет отвечать за счетчики.
У него будет следующая функциональность:
\begin{enumerate}
\item При конструировании заводит счетчики для всех шести методов и инициализирует их нулями.

\item Для каждого счетчика есть метод увеличения его значения на единицу вверх.

\item В деструкторе выводится информация пользователю о значениях счетчиков.
\end{enumerate}
Имплементировать класс можно следующим образом:
\begin{cppcode}
class Counter {
public:
  Counter() = default;
  ~Counter() {
    print_counters();
  }

  void increase_ctor() {
    ++ctor_counter_;
  }
  void increase_copy_ctor() {
    ++copy_ctor_counter_;
  }
  void increase_move_ctor() {
    ++move_ctor_counter_;
  }
  void increase_copy_asgmnt() {
    ++copy_asgmnt_counter_;
  }
  void increase_move_asgmnt() {
    ++move_asgmnt_counter_;
  }
  void increase_dtor() {
    ++dtor_counter_;
  }
  
private:
  void print_counter() const;

  uint64_t ctor_counter_ = 0;
  uint64_t copy_ctor_counter_ = 0;
  uint64_t move_ctor_counter_ = 0;
  uint64_t copy_asgmnt_counter_ = 0;
  uint64_t move_asgmnt_counter_ = 0;
  uint64_t dtor_counter_ = 0;
};
\end{cppcode}
Для красоты можно еще удалить копирующие и мувающие операции на счетчике.
Имплементация печати может быть такой:
\begin{cppcode}
void Counter::print_counter() const {
  std::cout << "Statistics:\n";
  std::cout << "ctors: " << ctor_counter_ << '\n';
  std::cout << "copys: " << copy_ctor_counter_ << '\n';
  std::cout << "moves: " << move_ctor_counter_ << '\n';
  std::cout << "copy=: " << copy_asgmnt_counter_ << '\n';
  std::cout << "move=: " << move_asgmnt_counter_ << '\n';
  std::cout << "dtors: " << dtor_counter_ << '\n';
}
\end{cppcode}
В этот класс можно добавить информацию об имени класса, для которого ведется подсчет, но я специально делаю минимальный работающий пример.

\paragraph{Имплементация \texttt{CTorCheck}}

Теперь имплементация класса \verb"CTorCheck", который будет выполнять вычисления со счетчиком:
\begin{cppcode}
template<class T>
class CTorCheck {
public:
  CTorCheck() {
    GCounter::instance().increase_ctor();
  }
  CTorCheck(const CTorCheck&) {
      GCounter::instance().increase_copy_ctor();
  }
  CTorCheck(CTorCheck&&) noexcept {
      GCounter::instance().increase_move_ctor();
  }
  CTorCheck& operator=(const CTorCheck&) {
    GCounter::instance().increase_copy_asgmnt();
    return *this;
  }
  CTorCheck& operator=(CTorCheck&&) noexcept {
    GCounter::instance().increase_move_asgmnt();
    return *this;
  }
  ~CTorCheck() {
    GCounter::instance().increase_dtor();
  }
private:
  using GCounter = Global<Counter, struct Tag>;
};
\end{cppcode}
Давайте поясним, что тут происходит.
\begin{enumerate}
\item В строке $25$ заводится глобальный объект типа \verb"Counter".
Обратите внимание, что полный тип будет следующий
\begin{cppcode}[numbers = none]
using GCounter = Global<Counter, CTorCheck<T>::Tag>;
\end{cppcode}
Таким образом для разных типов \verb"T" будут созданы разные счетчики.

\item В строках с $4$-й по $23$-ю каждый метод просто увеличивает свой счетчик внутри глобального \verb"GCounter".
\end{enumerate}
Такая имплементация создает глобальный счетчик при первом вызове.
А потому ее можно использовать даже для подсчета количества операций для глобальных объектов.
